% !TeX root = ../example-en.tex
\section{Media and closing}
\sectioncoverpage[\campusexamplesectionfiveimage]{Video captions, speaker notes, references and the closing page.}
\note{Five subsections demonstrate an odd number of entries in two columns. The last entry stays in the left column.}
\subsection{Video links}
\begin{frame}{Video example}
  \centering
  \href{run:../assets/presentation_demo.mp4}{%
    \XeTeXLinkBox{\parbox[c][.60\textheight][c]{\dimexpr .60\textheight*16/9\relax}{\mbox{}}}%
  }
  \campusvideocaption{From an outline to slides and speaker notes.}
\end{frame}
\note{The original six-second geometric animation shows outline items becoming slides and aligned notes. It is CC0-1.0, without external footage or audio. Clicking in the PDF opens the local video; PPTX export embeds it with a first-frame poster.}

\subsection{References and closing}
\begin{frame}[fragile]{Reference and closing commands}
  \begin{description}[More pages]
    \item[Single page] \texttt{\string\referencespage[shrink][\string\small]} sets a font and fits one page.
    \item[More pages] \texttt{\string\referencespage[break][\string\footnotesize]} allows continuation pages.
    \item[Columns] Append \texttt{[twocolumn]}; the default is one column.
    \item[Closing] \texttt{\string\backmatter} repeats the cover metadata; \texttt{[notitle]} hides the main title.
  \end{description}
\end{frame}
\note{Explain the single-page, continuation and two-column interfaces. The closing page repeats title, subtitle, course or group, author, advisor and date. The main title can be hidden.}

\subsection{Speaker notes}
\begin{frame}[fragile,label=demo-notes]{Slides and speaker notes}
  \begin{columns}[T,onlytextwidth]
    \begin{column}{.48\textwidth}
      \textbf{Source for one page}
      \begin{verbatim}
\begin{itemize}
  \item Main point
\end{itemize}
\note{Explain and connect.}
      \end{verbatim}
      Keep the slide and its notes in the same TeX source.
    \end{column}
    \begin{column}{.48\textwidth}
      \begin{description}[Notes]
        \item[PDF] \texttt{shownotes=false} shows slides only.
        \item[Notes] \texttt{shownotes=true} places notes on the right.
        \item[PPTX] TeX notes export to PowerPoint speaker notes.
      \end{description}
    \end{column}
  \end{columns}
\end{frame}
\note{Keep the main point on the slide and explanations, transitions and speaking cues in note. The shownotes class option defaults to false. Enabling it produces a slide-plus-notes PDF for a compatible viewer. PPTX always exports notes from the same TeX source, independently of whether the PDF displays them. Use note<1> and note<2> for stage-specific prompts.}

\subsection{Bibliography}
\referencespage[break][\small][onecolumn]
\note{The list collects the two cited entries in order of first citation. It demonstrates the full bibliography rather than the custom corner prefix.}

\subsection{Closing}
\backmatter
\note{Close the tour and invite feedback on code windows, readability and the authoring workflow.}
